% GENERATED FILE. Edit canonical lessons, then run npm run generate:books.
% canonical-source-hash: fnv1a64:c254e47955aea60d
% canonical-lessons: AR-C65-takallama, AR-C65-hamala, AR-C65-faza, AR-C65-nazala, AR-C65-samak

\chapter{To speak, To carry, To win, To go down, Fish}
\label{ch:ar-to-speak-to-carry-to-win-to-go-down-fish}
\hlchaptermodality{65}

\begin{chapteropening}
\textbf{By the end of this chapter:} \emph{I can say to speak, to carry, to win, to go down and fish.}

Complete the last lesson of chapter 65: I can say to speak, to carry, to win, to go down and fish.
\end{chapteropening}

\section[\texorpdfstring{\ar{تكلم} (takallama)}{تكلم (takallama)}]{\ar{تكلم} (takallama) — to speak (he spoke)}

\label{lesson:AR-C65-takallama}



\begin{quote}
Before the new one: say the Arabic for to wait, then the Arabic for to come back.
\end{quote}

\subsection*{You'll want to know: \ar{تكلم}}
\textbf{\ar{تكلم}} — \emph{takallama} — \textquotedblleft{}to speak (he spoke)\textquotedblright{}.

\begin{cousinweb}
\textbf{\ar{تكلم}} (\emph{takallama}) is built on the native Arabic root \emph{k-l-m}, \textquotedblleft{}to speak to; speech\textquotedblright{} (Lane). \emph{Takallama} is Form V, \textquotedblleft{}to speak\textquotedblright{}; the same root gives \emph{kalima}, \textquotedblleft{}a word\textquotedblright{}, and \emph{kalām}, \textquotedblleft{}speech\textquotedblright{}.
\end{cousinweb}

\begin{grammarlens}[title={What you've built}]
One more word for this chapter.
\end{grammarlens}

\subsection*{Your turn}
\begin{itemize}
\raggedright
  \item \emph{Say it:} \emph{takallama}
  \item \emph{Say it:} \emph{takallama}, once more
  \item \emph{Say it:} say \emph{takallama}
  \item \emph{Recall:} say the Arabic for to wait, then the Arabic for to come back, then say \emph{takallama} again
\end{itemize}

\begin{tcolorbox}[breakable,colback=teal!4,colframe=teal!35!black,title={Before you move on}]
What does \ar{تكلم} mean? (\textquotedblleft{}To speak (he spoke)\textquotedblright{}.) Say it once more.
\end{tcolorbox}
\section[\texorpdfstring{\ar{حمل} (ḥamala)}{حمل (ḥamala)}]{\ar{حمل} (ḥamala) — to carry (he carried)}

\label{lesson:AR-C65-hamala}



\begin{quote}
Before the new one: say the Arabic for to come back, then the Arabic for to speak.
\end{quote}

\subsection*{You'll want to know: \ar{حمل}}
\textbf{\ar{حمل}} — \emph{ḥamala} — \textquotedblleft{}to carry (he carried)\textquotedblright{}.

\begin{cousinweb}
\textbf{\ar{حمل}} (\emph{ḥamala}) is built on the native Arabic root \emph{ḥ-m-l}, \textquotedblleft{}to carry, to bear\textquotedblright{} (Lane). \emph{Ḥamala} is the plain verb, \textquotedblleft{}he carried\textquotedblright{}.
\end{cousinweb}

\begin{grammarlens}[title={What you've built}]
One more word for this chapter.
\end{grammarlens}

\subsection*{Your turn}
\begin{itemize}
\raggedright
  \item \emph{Say it:} \emph{ḥamala}
  \item \emph{Say it:} \emph{ḥamala}, once more
  \item \emph{Say it:} say \emph{ḥamala}
  \item \emph{Recall:} say the Arabic for to come back, then the Arabic for to speak, then say \emph{ḥamala} again
\end{itemize}

\begin{tcolorbox}[breakable,colback=teal!4,colframe=teal!35!black,title={Before you move on}]
What does \ar{حمل} mean? (\textquotedblleft{}To carry (he carried)\textquotedblright{}.) Say it once more.
\end{tcolorbox}
\section[\texorpdfstring{\ar{فاز} (fāza)}{فاز (fāza)}]{\ar{فاز} (fāza) — to win (he won)}

\label{lesson:AR-C65-faza}



\begin{quote}
Before the new one: say the Arabic for to speak, then the Arabic for to carry.
\end{quote}

\subsection*{You'll want to know: \ar{فاز}}
\textbf{\ar{فاز}} — \emph{fāza} — \textquotedblleft{}to win (he won)\textquotedblright{}.

\begin{cousinweb}
\textbf{\ar{فاز}} (\emph{fāza}) is built on the native Arabic root \emph{f-w-z}, \textquotedblleft{}to attain, to succeed\textquotedblright{} (Lane). \emph{Fāza} is the plain verb, \textquotedblleft{}he won\textquotedblright{}.
\end{cousinweb}

\begin{grammarlens}[title={What you've built}]
One more word for this chapter.
\end{grammarlens}

\subsection*{Your turn}
\begin{itemize}
\raggedright
  \item \emph{Say it:} \emph{fāza}
  \item \emph{Say it:} \emph{fāza}, once more
  \item \emph{Say it:} say \emph{fāza}
  \item \emph{Recall:} say the Arabic for to speak, then the Arabic for to carry, then say \emph{fāza} again
\end{itemize}

\begin{tcolorbox}[breakable,colback=teal!4,colframe=teal!35!black,title={Before you move on}]
What does \ar{فاز} mean? (\textquotedblleft{}To win (he won)\textquotedblright{}.) Say it once more.
\end{tcolorbox}
\section[\texorpdfstring{\ar{نزل} (nazala)}{نزل (nazala)}]{\ar{نزل} (nazala) — to go down (he went down)}

\label{lesson:AR-C65-nazala}



\begin{quote}
Before the new one: say the Arabic for to carry, then the Arabic for to win.
\end{quote}

\subsection*{You'll want to know: \ar{نزل}}
\textbf{\ar{نزل}} — \emph{nazala} — \textquotedblleft{}to go down (he went down)\textquotedblright{}.

\begin{cousinweb}
\textbf{\ar{نزل}} (\emph{nazala}) is built on the native Arabic root \emph{n-z-l}, \textquotedblleft{}to come down, to alight\textquotedblright{} (Lane). The same root gives \emph{manzil}, \textquotedblleft{}a place of alighting, a house\textquotedblright{}.
\end{cousinweb}

\begin{grammarlens}[title={What you've built}]
One more word for this chapter.
\end{grammarlens}

\subsection*{Your turn}
\begin{itemize}
\raggedright
  \item \emph{Say it:} \emph{nazala}
  \item \emph{Say it:} \emph{nazala}, once more
  \item \emph{Say it:} say \emph{nazala}
  \item \emph{Recall:} say the Arabic for to carry, then the Arabic for to win, then say \emph{nazala} again
\end{itemize}

\begin{tcolorbox}[breakable,colback=teal!4,colframe=teal!35!black,title={Before you move on}]
What does \ar{نزل} mean? (\textquotedblleft{}To go down (he went down)\textquotedblright{}.) Say it once more.
\end{tcolorbox}
\section[\texorpdfstring{\ar{سمك} (samak)}{سمك (samak)}]{\ar{سمك} (samak) — fish}

\label{lesson:AR-C65-samak}



\begin{quote}
Before the new one: say the Arabic for to win, then the Arabic for to go down.
\end{quote}

\subsection*{You'll want to know: \ar{سمك}}
\textbf{\ar{سمك}} — \emph{samak} — \textquotedblleft{}fish\textquotedblright{}.

\begin{cousinweb}
\textbf{\ar{سمك}} (\emph{samak}), \textquotedblleft{}fish\textquotedblright{}, is an old native noun, filed under the root \emph{s-m-k} (Lane), beside the verb \emph{samaka}, \textquotedblleft{}to rise, to be high\textquotedblright{}.
\end{cousinweb}

\begin{grammarlens}[title={What you've built}]
One more word for this chapter.
\end{grammarlens}

\subsection*{Your turn}
\begin{itemize}
\raggedright
  \item \emph{Say it:} \emph{samak}
  \item \emph{Say it:} \emph{samak}, once more
  \item \emph{Say it:} say \emph{samak}
  \item \emph{Recall:} say the Arabic for to win, then the Arabic for to go down, then say \emph{samak} again
\end{itemize}

\begin{tcolorbox}[breakable,colback=teal!4,colframe=teal!35!black,title={Before you move on}]
What does \ar{سمك} mean? (\textquotedblleft{}Fish\textquotedblright{}.) Say it once more.
\end{tcolorbox}
